% Source: https://artofproblemsolving.com/community/c3h2427111_several_problems_that_i_made
% Stable identity: aops-topic-2427111-q26
Let $f(t)=5\sqrt{5+t}$, so $x_{k+1}=f(x_k)$, and $f$ is increasing. Let $L=\dfrac{25+15\sqrt5}{2}$, the positive root of $L^2-25L-125=0$; then $L^2=25(5+L)$ and $L>0$, so $f(L)=L$.

\textbf{Convergence.} $x_1=5\sqrt5\approx11.2<L\approx29.3$. If $x_k<L$, then $x_{k+1}=f(x_k)<f(L)=L$; so $x_k<L$ for all $k$. Also $x_2=5\sqrt{5+5\sqrt5}>5\sqrt5=x_1$, and if $x_{k+1}>x_k$ then $x_{k+2}=f(x_{k+1})>f(x_k)=x_{k+1}$. So the sequence is increasing and bounded above, hence it converges to some $\ell$ with $0<\ell\le L$.

\textbf{The limit.} Since $f$ is continuous, $\ell=f(\ell)=5\sqrt{5+\ell}$, so $\ell^2=125+25\ell$, i.e. $\ell^2-25\ell-125=0$, whose roots are $\dfrac{25\pm15\sqrt5}{2}$. The root $\dfrac{25-15\sqrt5}{2}$ is negative, so
\[\ell=\boxed{\dfrac{25+15\sqrt5}{2}}\approx29.27.\]
